\section{Métrica de Kerr}

La primera vez que se menciona la métrica de Kerr es en el artículo de Kerr \cite{kerr1963gravitational}, en el año de 1963. Esta métrica corresponde a la solución de la ecuación de campo de Einstein para un agujero negro en rotación.

Originalmente, la métrica se expresó para una geodésica nula (véase \ref{def:null_geodesic}), mediante el sistema de coordenadas Eddington-Finkelstein \cite{finkelstein1958past}, de la siguiente manera:

\begin{align}
ds^2 = &(r^2 + a^2 \cos^2\theta)\left(d\theta^2 + \sin^2\theta \, d\phi^2\right) \\
& + 2(du + a \sin^2\theta \, d\phi)\left(dr + a \sin^2\theta \, d\phi\right) \\
&- \left(1 - \frac{2mr}{r^2 + a^2 \cos^2\theta}\right)
\left(du + a \sin^2\theta \, d\phi\right)^2 \label{eq:metrica_de_kerr}
\end{align}

Cabe aclarar que las coordenadas no corresponden a las coordenadas esféricas.

El sistema de coordenadas de la ecuación \eqref{eq:metrica_de_kerr} no es el único con el que puede expresarse la métrica de Kerr. Una forma alternativa, adecuada para el estudio de las geodésicas de tipo tiempo (véase \ref{def:timelike_geodesic}), todavía contiene términos cruzados en la coordenada radial, del tipo $dr\,d\varphi'$ y $dr\,dt'$:

\begin{align}
    ds^2 = & \ dr^2 - 2a \sin^2\theta \, dr \, d\varphi' \\
    & + (r^2 + a^2)\sin^2\theta \, d\varphi'^2 \\
    & + \Sigma \, d\theta^2 - dt'^2 \\
    & + \frac{2mr}{\Sigma} \left( dr - a \sin^2\theta \, d\varphi' - dt' \right)^2
\end{align}

Boyer y Lindquist \cite{boyer1967maximal} introdujeron una transformación de coordenadas que elimina precisamente estos términos cruzados en $dr$. Las coordenadas de Boyer-Lindquist $(\bar{t}, \bar{r}, \bar{\theta}, \bar{\varphi})$ se definen a partir de las coordenadas primadas mediante:

\begin{align}
    \bar{r} &= r', \qquad \bar{\theta} = \theta', \\
    d\bar{\varphi} &= d\varphi' - \frac{a}{\Delta}\, dr', \\
    d\bar{t} &= dt' - \frac{2mr'}{\Delta}\, dr',
\end{align}

con $\Delta = r^2 - 2mr + a^2$. Bajo esta transformación la métrica pierde los términos cruzados con $dr$ y queda diagonal salvo por el único término cruzado $t$-$\varphi$ que presentamos a continuación. Estas coordenadas tampoco corresponden a las coordenadas esféricas.

Como se observa, el tipo de coordenadas que se usa y su transformación es importante, no es un detalle menor. Sin embargo, en la práctica manejar todas estas distinciones consume mucho tiempo y vuelve complejo el sistema. Por ello, asumiremos que en la respectiva métrica con la que estamos trabajando se entiende que se están usando las coordenadas que simplifican la expresión de la métrica. En el caso de la métrica de Kerr, a no ser que se especifique lo contrario, se asumirá que se están usando las coordenadas de Boyer-Lindquist. Por lo tanto, de ahora en adelante no primaremos las coordenadas; en su lugar escribiremos de la forma $x^\nu$, de modo que la conversión a cada uno de los sistemas que conocemos será de la forma $x^\nu = f(x'^\mu)$.

\subsubsection{Métrica de Kerr en coordenadas de Boyer-Lindquist}
La forma en la que, por simplicidad, decidimos usar la métrica de Kerr en coordenadas de Boyer-Lindquist es la siguiente:

\begin{align}
    ds^2 = A (dx^0)^2 + B (dx^1)^2 + C (dx^2)^2 + D (dx^3)^2 + E (dx^0)(dx^3)
\end{align}

Con:

\begin{align}
    A &= 1 - R_s  x^1 / \Sigma\\
    B &= -\Sigma / \Delta \\
    C &= -\Sigma \\
    D &= -\left((x^1)^2 + a^2 + R_s  x^1  \frac{a^2}{ \Sigma}  \sin^2(x^2)\right)  \sin^2(x^2) \\
    E &= 2  R_s  x^1  \frac{a}{ \Sigma}  \sin^2(x^2)
\end{align}

donde:
\begin{align}
    \Sigma &= r^2 + a^2 \cos^2(x^2) \\
    \Delta &= r^2 - 2mr + a^2 \\
    R_s &= 2GM/c^2 \\
    a &= \frac{J}{M}
\end{align}

Esta forma de la métrica está basada en el artículo de Teukolsky \cite{teukolsky2015kerr}, pero introduciendo el radio de Schwarzschild.


\subsubsection{Solución analítica de las geodésicas de tipo tiempo en la métrica de Kerr}
En 2009, Fujita y Hikida \cite{fujita2009analytical} encontraron una solución analítica para las geodésicas en la métrica de Kerr, basado en el tiempo de Mino \cite{mino2003perturbative}, una parametrización basada en las coordenadas de Boyer-Lindquist, que permite obtener una solución en términos de integrales elípticas de Legendre de primera, segunda y tercera especie, e integrales elípticas de Jacobi.

Las geodésicas tipo tiempo en la métrica de Kerr admiten tres constantes de movimiento por unidad de masa: la energía específica $E$, la componente $z$ del momento angular específico $L_z$, y la constante de Carter $Q$ \cite{carter1968global}. Estas constantes surgen de las simetrías de la métrica y permiten separar completamente las ecuaciones de movimiento.

Las ecuaciones de movimiento en tiempo propio $\tau$, en coordenadas de Boyer-Lindquist, toman la forma:

\begin{align}
    \Sigma^2 \left(\frac{dr}{d\tau}\right)^2 &= R(r), \label{eq:geodesica_r_tau} \\
    \Sigma^2 \left(\frac{d\cos\theta}{d\tau}\right)^2 &= \Theta(\cos\theta), \label{eq:geodesica_theta_tau} \\
    \Sigma \frac{dt}{d\tau} &= T_r(r) + T_\theta(\cos\theta), \label{eq:geodesica_t_tau} \\
    \Sigma \frac{d\varphi}{d\tau} &= \Phi_r(r) + \Phi_\theta(\cos\theta), \label{eq:geodesica_phi_tau}
\end{align}

En las ecuaciones \eqref{eq:geodesica_r_tau}, \eqref{eq:geodesica_theta_tau}, \eqref{eq:geodesica_t_tau} y \eqref{eq:geodesica_phi_tau}, el acoplamiento a través de $\Sigma = r^2 + a^2\cos^2\theta$ impide separar directamente las ecuaciones radial y polar.

El tiempo de Mino $\lambda$ se define mediante la reparametrización \cite{mino2003perturbative}:

\begin{align}
    d\tau = \Sigma \, d\lambda, \qquad \Sigma = r^2 + a^2\cos^2\theta. \label{eq:tiempo_de_mino}
\end{align}

Con la reparametrización \eqref{eq:tiempo_de_mino}, las ecuaciones de movimiento toman la forma \cite{fujita2009analytical}:

\begin{align}
    \left(\frac{dr}{d\lambda}\right)^2 &= R(r), \label{eq:geodesica_r_mino} \\
    \left(\frac{d\cos\theta}{d\lambda}\right)^2 &= \Theta(\cos\theta), \label{eq:geodesica_theta_mino} \\
    \frac{dt}{d\lambda} &= T_r(r) + T_\theta(\cos\theta), \label{eq:geodesica_t_mino} \\
    \frac{d\varphi}{d\lambda} &= \Phi_r(r) + \Phi_\theta(\cos\theta). \label{eq:geodesica_phi_mino}
\end{align}

La propiedad central del tiempo de Mino es que las ecuaciones \eqref{eq:geodesica_r_mino} y \eqref{eq:geodesica_theta_mino} quedan completamente desacopladas: la primera depende únicamente de $r$ y la segunda únicamente de $\theta$. Definiendo $P(r) = E(r^2 + a^2) - aL_z$, los potenciales y funciones auxiliares son:

\begin{align}
    R(r) &= \left[P(r)\right]^2 - \Delta\left[r^2 + (aE - L_z)^2 + Q\right], \label{eq:potencial_R} \\
    \Theta(\cos\theta) &= Q - \left(Q + a^2(1-E^2) + L_z^2\right)\cos^2\theta + a^2(1-E^2)\cos^4\theta, \label{eq:potencial_Theta} \\
    T_r(r) &= \frac{r^2 + a^2}{\Delta}\,P(r), \qquad T_\theta(\cos\theta) = -a^2 E\sin^2\theta + aL_z, \\
    \Phi_r(r) &= \frac{a}{\Delta}\,P(r), \qquad \Phi_\theta(\cos\theta) = \frac{L_z}{\sin^2\theta} - aE.
\end{align}

Para órbitas acotadas, el potencial radial $R(r)$ de \eqref{eq:potencial_R} posee cuatro raíces reales $r_1 \geq r_2 \geq r_3 \geq r_4$, donde $r_1$ es el apoápside y $r_2$ el periápside. El potencial polar $\Theta(\cos\theta)$ de \eqref{eq:potencial_Theta} se factoriza como:

\begin{align}
    \Theta(\cos\theta) = L_z^2\,\varepsilon_0\left(z_- - \cos^2\theta\right)\left(z_+ - \cos^2\theta\right), \label{eq:factorizacion_Theta}
\end{align}

En la factorización \eqref{eq:factorizacion_Theta}, $\varepsilon_0 = a^2(1-E^2)/L_z^2$, y $z_-$, $z_+$ son las raíces de $\Theta$ en la variable $\cos^2\theta$. La reducción a integrales elípticas se realiza mediante las transformaciones de variable y los módulos:

\begin{align}
    y_r &= \sqrt{\frac{r_1 - r_3}{r_1 - r_2}\cdot\frac{r - r_2}{r - r_3}}, \qquad
    k_r = \sqrt{\frac{r_1 - r_2}{r_1 - r_3}\cdot\frac{r_3 - r_4}{r_2 - r_4}}, \label{eq:modulo_kr} \\
    y_\theta &= \frac{\cos\theta}{\sqrt{z_-}}, \qquad
    k_\theta = \sqrt{\frac{z_-}{z_+}}. \label{eq:modulo_ktheta}
\end{align}

Con los módulos \eqref{eq:modulo_kr} y \eqref{eq:modulo_ktheta}, las integrales de camino desde los puntos de retorno se reducen a integrales elípticas incompletas de primera especie $F(\varphi, k)$:

\begin{align}
    \int_{r_2}^{r} \frac{dr'}{\sqrt{R(r')}} &= \frac{2}{\sqrt{(1-E^2)(r_1-r_3)(r_2-r_4)}}\,F\!\left(\arcsin y_r,\, k_r\right), \label{eq:integral_r} \\
    \int_{0}^{\cos\theta} \frac{d\cos\theta'}{\sqrt{\Theta(\cos\theta')}} &= \frac{1}{L_z\sqrt{\varepsilon_0 z_+}}\,F\!\left(\arcsin y_\theta,\, k_\theta\right). \label{eq:integral_theta}
\end{align}

Los períodos en tiempo de Mino se obtienen evaluando las integrales \eqref{eq:integral_r} y \eqref{eq:integral_theta} en los puntos de retorno:

\begin{align}
    \Lambda_r &= \frac{4\,K(k_r)}{\sqrt{(1-E^2)(r_1-r_3)(r_2-r_4)}}, \qquad
    \Lambda_\theta = \frac{4\,K(k_\theta)}{L_z\sqrt{\varepsilon_0 z_+}}, \label{eq:periodos_mino}
\end{align}

donde $K(k) = F(\pi/2, k)$ es la integral elíptica completa de primera especie; $\Lambda_r$ y $\Lambda_\theta$ en \eqref{eq:periodos_mino} corresponden a un período completo de oscilación en $r$ y en $\theta$, respectivamente.

La solución analítica para la coordenada radial, en términos de la función seno elíptico de Jacobi $\mathrm{sn}(u,k)$, es \cite{fujita2009analytical}:

\begin{align}
    r(\lambda) = \frac{r_3(r_1 - r_2)\,\mathrm{sn}^2(u_r,\,k_r) - r_2(r_1 - r_3)}{(r_1 - r_2)\,\mathrm{sn}^2(u_r,\,k_r) - (r_1 - r_3)}, \label{eq:solucion_r}
\end{align}

y para la coordenada polar:

\begin{align}
    \cos\theta(\lambda) = \sqrt{z_-}\,\mathrm{sn}(u_\theta,\,k_\theta), \label{eq:solucion_theta}
\end{align}

En las soluciones \eqref{eq:solucion_r} y \eqref{eq:solucion_theta}, $u_r$ y $u_\theta$ son funciones lineales a trozos de $\lambda$ que rastrean la fase dentro de cada período.

Las coordenadas $t(\lambda)$ y $\varphi(\lambda)$ se obtienen integrando \eqref{eq:geodesica_t_mino} y \eqref{eq:geodesica_phi_mino}. La solución se descompone en una parte secular y una parte oscilatoria \cite{fujita2009analytical}:

\begin{align}
    t(\lambda) &= \Gamma\lambda + t^{(r)}(\lambda) + t^{(\theta)}(\lambda), \label{eq:solucion_t} \\
    \varphi(\lambda) &= \Upsilon_\varphi\lambda + \varphi^{(r)}(\lambda) + \varphi^{(\theta)}(\lambda), \label{eq:solucion_phi}
\end{align}

En las soluciones \eqref{eq:solucion_t} y \eqref{eq:solucion_phi}, $\Gamma$ y $\Upsilon_\varphi$ son las frecuencias medias en tiempo de Mino, expresadas en términos de integrales elípticas completas de primera, segunda y tercera especie. Las partes oscilatorias $t^{(r)}$, $t^{(\theta)}$, $\varphi^{(r)}$, $\varphi^{(\theta)}$ resultan en combinaciones de integrales elípticas incompletas de segunda y tercera especie, evaluadas a lo largo de la órbita. Las frecuencias físicas observables \eqref{eq:frecuencias_fisicas} se relacionan con las frecuencias en tiempo de Mino mediante:

\begin{align}
    \Omega_r = \frac{\Upsilon_r}{\Gamma}, \qquad \Omega_\theta = \frac{\Upsilon_\theta}{\Gamma}, \qquad \Omega_\varphi = \frac{\Upsilon_\varphi}{\Gamma}, \label{eq:frecuencias_fisicas}
\end{align}

donde $\Upsilon_r = 2\pi/\Lambda_r$ y $\Upsilon_\theta = 2\pi/\Lambda_\theta$ son las frecuencias angulares en tiempo de Mino.

Estas soluciones analíticas presuponen las constantes de movimiento $E$, $L_z$ y $Q$ que aparecen en los potenciales; en la siguiente sección precisamos su definición y las ecuaciones de la geodésica de las que se derivan.

